\chapter{Complexity — Gate}
%%% generated by the blueprint pipeline from TCSlib.Complexity.FormulaKit.Gate
\section{Overview}

Declarations of the Lean module \texttt{TCSlib.Complexity.FormulaKit.Gate}.

\section{Declarations}

\begin{definition}[Block variable index]
\lean{Complexity.FormulaKit.blockVar}
\label{Complexity.FormulaKit.blockVar}
\leanok
For natural numbers $b,w,t,i$, the block variable index is $b+tw+i$.
\end{definition}

\begin{definition}[Bits of a natural number]
\lean{Complexity.FormulaKit.bitsOf}
\label{Complexity.FormulaKit.bitsOf}
\leanok
For natural numbers $ell$ and $r$, the function on indices $i<ell$ whose value at $i$
is the $i$th binary digit of $r$.
\end{definition}

\begin{definition}[Clause excluding an assignment]
\lean{Complexity.FormulaKit.excludeClause}
\label{Complexity.FormulaKit.excludeClause}
\leanok
For a natural number $ell$ and a Boolean assignment $v$ on the indices
$0,dots,ell-1$, form the clause containing, in index order, the literal on variable
$i$ with Boolean value opposite to $v(i)$.
\end{definition}

\begin{definition}[Boolean function template formula]
\lean{Complexity.FormulaKit.template}
\label{Complexity.FormulaKit.template}
\leanok
\uses{Complexity.FormulaKit.bitsOf, Complexity.FormulaKit.excludeClause}
For a natural number $ell$ and a Boolean function $p$ on assignments to the indices
below $ell$, form a CNF formula by considering each natural number $r<2^ell$ in
increasing order: include the clause excluding the assignment given by the binary digits
of $r$ exactly when $p$ is false on that assignment.
\end{definition}

\begin{theorem}[Evaluation of the template formula]
\lean{Complexity.FormulaKit.template_eval}
\label{Complexity.FormulaKit.template_eval}
\leanok
\uses{Complexity.FormulaKit.template}
For every natural number $\ell$, every Boolean predicate $p$ on strings of length
$\ell$, and every assignment $a:\mathbb{N}\to\{\mathrm{false},\mathrm{true}\}$, the CNF
template formed by excluding each length-$\ell$ Boolean string on which $p$ is false
evaluates under $a$ to $p$ applied to the restriction $i\mapsto a(i)$ for $i<\ell$.
\end{theorem}

\begin{theorem}[Template variable bound]
\lean{Complexity.FormulaKit.template_numVars_le}
\label{Complexity.FormulaKit.template_numVars_le}
\leanok
\uses{Complexity.FormulaKit.template, Std.Sat.CNF.numVars}
For every natural number $\ell$ and every Boolean predicate on strings of length $\ell$,
the variable count of its CNF template is at most $\ell$.
\end{theorem}

\begin{theorem}[Template length bound]
\lean{Complexity.FormulaKit.template_length_le}
\label{Complexity.FormulaKit.template_length_le}
\leanok
\uses{Complexity.FormulaKit.template}
For every natural number $\ell$ and every Boolean predicate on strings of length $\ell$,
the number of clauses in its CNF template is at most $2^\ell$.
\end{theorem}

\begin{theorem}[Template clause width]
\lean{Complexity.FormulaKit.template_widthAtMost}
\label{Complexity.FormulaKit.template_widthAtMost}
\leanok
\uses{Complexity.FormulaKit.template, Std.Sat.CNF.WidthAtMost}
For every natural number $\ell$ and every Boolean predicate on strings of length $\ell$,
every clause in its CNF template has at most $\ell$ literals; thus the template has
width at most $\ell$.
\end{theorem}

\begin{theorem}[Variable count bound from index bounds]
\lean{Complexity.FormulaKit.numVars_le_of_forall_lt}
\label{Complexity.FormulaKit.numVars_le_of_forall_lt}
\leanok
\uses{Std.Sat.CNF.numVars}
Let $\varphi$ be a CNF formula over natural-number variables, and let $N$ be a natural
number. If every variable index occurring in every clause of $\varphi$ is less than $N$,
then the variable count of $\varphi$ is at most $N$.
\end{theorem}

\begin{definition}[Relabelled Boolean gate formula]
\lean{Complexity.FormulaKit.gate}
\label{Complexity.FormulaKit.gate}
\leanok
\uses{Complexity.FormulaKit.template}
For a natural number $ell$, a Boolean function $p$ on assignments to the indices below
$ell$, and a map $w$ from natural-number variables to natural-number variables, the
gate formula is the template formula for $p$ with each variable relabelled by $w$.
\end{definition}

\begin{theorem}[Evaluation of a gate]
\lean{Complexity.FormulaKit.gate_eval}
\label{Complexity.FormulaKit.gate_eval}
\leanok
\uses{Complexity.FormulaKit.gate, Complexity.FormulaKit.template_eval}
\difficulty{3}
For every natural number $\ell$, every Boolean function $p$ on Boolean vectors indexed
by $\{0,\ldots,\ell-1\}$, every map $\mathrm{wire}:\mathbb{N}\to\mathbb{N}$, and every
Boolean assignment $a$ on the natural numbers, the gate evaluates to $p$ applied to the
vector whose value at index $i$ is $a(\mathrm{wire}(i))$.
\end{theorem}

\begin{definition}[Bit read with default value]
\lean{Complexity.FormulaKit.readBit}
\label{Complexity.FormulaKit.readBit}
\leanok
For a natural number $ell$, a Boolean assignment $v$ on indices below $ell$, and a
natural number $j$, read the value $v(j)$ when $j<ell$, and return false otherwise.
\end{definition}

\begin{definition}[Exactly-one constraint formula]
\lean{Complexity.FormulaKit.exactlyOne}
\label{Complexity.FormulaKit.exactlyOne}
\leanok
For a list $v_s$ of natural-number variables, form a CNF formula whose first clause
contains the positive literal for each variable in $v_s$, in list order. For each
variable $v$ in the list, also include a negative-literal clause for every variable $w$
in the list with $w\ne v$, consisting of the negative literals for $v$ and $w$ in that
order.
\end{definition}

\begin{theorem}[Evaluation of the exactly-one formula]
\lean{Complexity.FormulaKit.exactlyOne_eval}
\label{Complexity.FormulaKit.exactlyOne_eval}
\leanok
\uses{Complexity.FormulaKit.exactlyOne}
For every finite list $vs$ of natural numbers and every Boolean assignment $a$ to the
natural numbers, the exactly-one formula evaluates to true if and only if there is some
$j$ appearing in $vs$ such that, for every $v$ appearing in $vs$, $a(v)$ is true exactly
when $v=j$.
\end{theorem}

\begin{definition}[CNF formula cost]
\lean{Complexity.FormulaKit.cost}
\label{Complexity.FormulaKit.cost}
\leanok
For a CNF formula $\varphi$ over a type $\alpha$, define its cost to be the sum, over
all clauses $C$ in the formula, of $|C|+1$, where $|C|$ is the clause length.
\end{definition}

\begin{theorem}[Cost invariance under relabeling]
\lean{Complexity.FormulaKit.cost_relabel}
\label{Complexity.FormulaKit.cost_relabel}
\leanok
\uses{Complexity.FormulaKit.cost}
\difficulty{3}
For any types $\alpha$ and $\beta$, any CNF formula $\varphi$ over $\alpha$, and any map
$f: \alpha \to \beta$, relabeling the variables by $f$ preserves the cost:
$\operatorname{cost}(\varphi\text{ relabeled by }f)=\operatorname{cost}(\varphi)$.
\end{theorem}

\begin{theorem}[Cost of Concatenated Formulas]
\lean{Complexity.FormulaKit.cost_append}
\label{Complexity.FormulaKit.cost_append}
\leanok
\uses{Complexity.FormulaKit.cost}
\difficulty{1}
For any two conjunctive normal forms $\varphi$ and $\psi$ over the same variable type,
the cost of their concatenation is the sum of their costs.
\end{theorem}

\begin{theorem}[Cost Bound for Flattened Formulas]
\lean{Complexity.FormulaKit.cost_flatMap_le}
\label{Complexity.FormulaKit.cost_flatMap_le}
\leanok
\uses{Complexity.FormulaKit.cost}
Let $xs$ be a finite list and assign to each element $x$ a conjunctive normal form
$f(x)$. If every $x$ occurring in $xs$ has cost at most $B$, then the cost of
concatenating all the formulas $f(x)$ for $x$ in $xs$ is at most $|xs|B$.
\end{theorem}

\begin{theorem}[Length bound for formula encoding]
\lean{Complexity.FormulaKit.serialize_length_le}
\label{Complexity.FormulaKit.serialize_length_le}
\leanok
\uses{Complexity.FormulaKit.cost, Std.Sat.CNF.serialize}
Let $\varphi$ be a CNF formula over variables indexed by the natural numbers, and let
$N\in\mathbb N$. Suppose every literal in every clause of $\varphi$ has variable index
less than $N$. If the cost of $\varphi$ is the sum, over its clauses $C$, of $|C|+1$,
then the length of its binary encoding is at most $\operatorname{cost}(\varphi)(N+3)+1$.
\end{theorem}

\begin{theorem}[Gate cost bound]
\lean{Complexity.FormulaKit.gate_cost_le}
\label{Complexity.FormulaKit.gate_cost_le}
\leanok
\uses{Complexity.FormulaKit.cost, Complexity.FormulaKit.gate}
For every natural number $\ell$, every Boolean function $p$ on Boolean vectors indexed
by $\{0,\ldots,\ell-1\}$, and every map $\mathrm{wire}:\mathbb{N}\to\mathbb{N}$, the
cost of the CNF gate obtained by relabeling the truth-table template for $p$ along
$\mathrm{wire}$ is at most $2^\ell(\ell+1)$.
\end{theorem}

\begin{theorem}[Cost bound for the exactly-one formula]
\lean{Complexity.FormulaKit.exactlyOne_cost_le}
\label{Complexity.FormulaKit.exactlyOne_cost_le}
\leanok
\uses{Complexity.FormulaKit.cost, Complexity.FormulaKit.exactlyOne}
For every finite list $vs$ of natural numbers, the cost of the CNF formula expressing
exactly one of the listed variables is at most $|vs|+1+3|vs|^2$, where $|vs|$ is the
length of the list.
\end{theorem}
